\documentclass[aps,prc,onecolumn,superscriptaddress,nofootinbib,10pt]{revtex4-2}
\usepackage{amsmath,amssymb}
\usepackage{booktabs}
\usepackage[colorlinks=true,linkcolor=blue,citecolor=blue,urlcolor=blue]{hyperref}

\newcommand{\nB}{n_B}
\newcommand{\nC}{n_C}
\newcommand{\nS}{n_S}
\newcommand{\muB}{\mu_B}
\newcommand{\muC}{\mu_C}
\newcommand{\muS}{\mu_S}
\newcommand{\munue}{\mu_{\nu_e}}
\newcommand{\hc}{(\hbar c)^3}
\newcommand{\as}{\alpha_s}

\begin{document}

\title{$\alpha$Bag: the perturbatively corrected bag model, unpaired and
       colour--flavour locked \\
       --- as implemented in \texttt{eos/alphabag}}

\author{M.~Guerrini}
\affiliation{University of Ferrara}

\begin{abstract}
The \texttt{eos/alphabag} subpackage implements deconfined quark matter as a
gas of $u$, $d$ and $s$ quarks inside a bag, with the leading perturbative QCD
correction carried as a single coupling $\as$ multiplying the free-gas
pressure, in the form used by Fischer \textit{et al.}~\cite{Fischer2011}. The
strange quark is massive and its kinetic thermodynamics is the exact Fermi
gas; the light flavours are massless and close in elementary functions.
Thermal gluons are an optional sector with their own $\as$ correction. On top
of the unpaired phase the package carries a second phase, colour--flavour
locked (CFL) quark matter~\cite{AlfordRajagopalWilczek1999}, obtained by
adding the pairing term of Alford, Braby, Paris and
Reddy~\cite{AlfordBrabyParisReddy2005} to the same potential and closing the
system with flavour locking $n_u = n_d = n_s$ instead of with an equilibrium
condition. This note states the thermodynamic potential of both phases, the
parameters, the conserved-charge bookkeeping, the residual of every mode the
package closes, and every quantity a solved point returns. It is written from
the code, as its specification: every equation here is asserted by
\texttt{eos/alphabag/verify/run\_full\_check.py} or by \texttt{test/baseline}.
\end{abstract}

\maketitle

\section{The thermodynamic potential}
\label{sec:potential}

Uniform deconfined matter at flavour potentials $\mu_u, \mu_d, \mu_s$ and
temperature $T$ is a gas of quarks and antiquarks confined by a constant
vacuum pressure $B$. The pressure of the strongly-interacting sector is
%
\begin{equation}
P \;=\; \sum_{q = u,d,s} P_q(\mu_q, T, m_q, \as)
        \;+\; P_g(T, \as) \;-\; \frac{B}{\hc} ,
\label{eq:P}
\end{equation}
%
and, since the bag constant is a constant shifted out of the vacuum energy,
%
\begin{equation}
\varepsilon \;=\; \sum_q \varepsilon_q \;+\; \varepsilon_g
                  \;+\; \frac{B}{\hc} , \qquad
s \;=\; \sum_q s_q + s_g , \qquad
n_q \;=\; n_q(\mu_q, T, m_q, \as) .
\label{eq:eps}
\end{equation}
%
The bag enters $\varepsilon$ and $P$ with opposite signs and enters neither
$s$ nor any $n_q$, so it cancels out of $\varepsilon + P$ and the Euler
relation holds with no bag term in it,
%
\begin{equation}
\varepsilon + P \;=\; T s + \sum_q \mu_q n_q ,
\label{eq:euler}
\end{equation}
%
which is checked in \texttt{verify/} and holds to $4\times 10^{-16}$
relative. The free energy density is $f = \varepsilon - Ts = -P + \sum_q \mu_q
n_q$.

Two remarks on what Eq.~\eqref{eq:P} is not. It carries no vector field: the
repulsion that in \texttt{eos/vmit} comes from a $\omega$-like mean field is
here replaced by the perturbative correction, which is a multiplicative
softening rather than a stiffening, so the two quark models of this repository
stiffen the equation of state by different mechanisms and are not
reparametrisations of one another. And it carries no scalar field and no gap
equation: the quark masses $m_u, m_d, m_s$ are parameters, not solutions, so
the potential is explicit in $\mu$ and the only solving to be done is that of
the equilibrium conditions themselves.

\subsection{The perturbative correction}
\label{sec:alpha}

The correction is the $\mathcal{O}(\as)$ term of the free energy of a
relativistic quark--gluon gas~\cite{FreedmanMcLerran1977}, kept as two
multiplicative factors --- one on the quark thermal term, one on everything
carrying a chemical potential ---
%
\begin{equation}
c_q(\as) \;=\; 1 - \frac{2\as}{\pi} , \qquad
c_T(\as) \;=\; 1 - \frac{50\,\as}{21 \pi} , \qquad
c_g(\as) \;=\; 1 - \frac{15\,\as}{4 \pi} ,
\label{eq:cfactors}
\end{equation}
%
in the arrangement of Fischer \textit{et al.}~\cite{Fischer2011}. At the
shipped $\as = 0.3$ these are $c_q = 0.8090$, $c_T = 0.7726$ and
$c_g = 0.6419$. There is no running: $\as$ is a constant of the parameter set,
which is what makes it a parameter an inference run can vary.

\subsection{Parameters}

\begin{center}
\begin{tabular}{llrl}
\toprule
symbol & code & value & meaning \\
\midrule
$m_u$      & \texttt{m\_u}   & $0$ MeV      & up-quark mass, treated as massless \\
$m_d$      & \texttt{m\_d}   & $0$ MeV      & down-quark mass, treated as massless \\
$m_s$      & \texttt{m\_s}   & $150$ MeV    & strange-quark mass \\
$\as$      & \texttt{alpha}  & $0.3$        & QCD coupling, constant \\
$B^{1/4}$  & \texttt{B4}     & $165$ MeV    & bag constant, quoted as its fourth root \\
$B$        & \texttt{B}      & $7.412\times 10^{8}\ \mathrm{MeV}^4$
                                            & $= (B^{1/4})^4 = 96.466\ \mathrm{MeV/fm}^3$ \\
\bottomrule
\end{tabular}
\end{center}
%
$B$ is a derived property of the parameter record, never stored, so a set
cannot carry a $B$ and a $B^{1/4}$ that disagree. The defaults are a central
choice within the ranges used for compact-star quark matter in
Refs.~\cite{AlfordReddy2003,Fischer2011}
($B^{1/4} \sim 145$--$180$~MeV, $\as \sim 0.2$--$0.5$,
$m_s \sim 90$--$150$~MeV); they are named defaults and not hardcoded values,
and every entry point takes the parameter record as its first argument.

The pairing gap $\Delta_0$ of the CFL phase (Sec.~\ref{sec:cfl}) is
deliberately \emph{not} a field of the parameter record. It selects a phase
rather than tuning one, in the same way the species flags do, and it is passed
per call.

For orientation, the shipped set at $T = 0$ in charge-neutral beta equilibrium
reaches $P = 0$ --- the surface of a self-bound star --- at
$\nB = 0.40309\ \mathrm{fm}^{-3}$ with $\varepsilon/\nB = 1046.17$~MeV, so
this set is not absolutely stable strange quark matter. Pairing lowers it:
the CFL phase at $\Delta_0 = 100$~MeV has its surface at
$\nB = 0.36286\ \mathrm{fm}^{-3}$ with $\varepsilon/\nB = 936.56$~MeV.

\section{Single-flavour thermodynamics}
\label{sec:kinetic}

Each flavour is a gas of quarks and antiquarks of degeneracy
$g = 2 \times 3 = 6$ (spin $\times$ colour). Two cases are implemented and the
choice is made on the mass, $m < 10^{-5}$~MeV selecting the massless branch.

\subsection{Massless flavours}

For $m = 0$ the Fermi integrals close in elementary functions, and with the
factors of Eq.~\eqref{eq:cfactors} folded in the flavour carries
%
\begin{align}
P_0(\mu, T, \as) &= \frac{1}{\hc} \left[
    \frac{7}{60}\pi^2 T^4\, c_T(\as)
    + \left( \frac{T^2 \mu^2}{2} + \frac{\mu^4}{4\pi^2} \right) c_q(\as)
    \right] ,
\label{eq:P0} \\
n_0(\mu, T, \as) &= \frac{1}{\hc}
    \left( \mu T^2 + \frac{\mu^3}{\pi^2} \right) c_q(\as) ,
\label{eq:n0} \\
s_0(\mu, T, \as) &= \frac{1}{\hc} \left[
    \frac{7}{15}\pi^2 T^3\, c_T(\as) + T \mu^2\, c_q(\as) \right] ,
\label{eq:s0} \\
\varepsilon_0(\mu, T, \as) &= 3\, P_0(\mu, T, \as) .
\label{eq:e0}
\end{align}
%
At $\as = 0$ these are the textbook massless Fermi gas at degeneracy $6$ with
antiparticles. They are a consistent set: $n_0 = \partial P_0/\partial\mu$ and
$s_0 = \partial P_0 / \partial T$ hold identically (verified against numerical
derivatives to $10^{-10}$ relative), and $\varepsilon_0 + P_0 = 4P_0 = T s_0 +
\mu n_0$ term by term, so Eq.~\eqref{eq:euler} holds for a massless flavour at
any $\as$. This is why the correction is applied to $P$, $n$, $\varepsilon$
and $s$ separately rather than to $P$ alone: applying it to $P$ and taking the
others as derivatives would give the same answer, and the code writes all four
out for readability.

At $T = 0$ Eqs.~\eqref{eq:P0}--\eqref{eq:e0} reduce to
$P_0 = \mu^4 c_q / (4\pi^2\hc)$, $n_0 = \mu^3 c_q/(\pi^2 \hc)$, $s_0 = 0$.

\subsection{Massive flavours}

For $m > 0$ the kinetic thermodynamics is the exact Fermi gas,
%
\begin{align}
n_F &= \frac{g}{2\pi^2 \hc} \int_0^\infty \! dk\, k^2
       \left[ f(E_k - \mu) - f(E_k + \mu) \right] ,
       & E_k &= \sqrt{k^2 + m^2} ,
\label{eq:nF} \\
P_F &= \frac{g}{6\pi^2 \hc} \int_0^\infty \! dk\, \frac{k^4}{E_k}
       \left[ f(E_k - \mu) + f(E_k + \mu) \right] ,
\label{eq:PF} \\
\varepsilon_F &= \frac{g}{2\pi^2 \hc} \int_0^\infty \! dk\, k^2 E_k
       \left[ f(E_k - \mu) + f(E_k + \mu) \right] ,
\label{eq:eF} \\
s_F &= \frac{\varepsilon_F + P_F - \mu\, n_F}{T} ,
\label{eq:sF}
\end{align}
%
with $f(x) = [1 + e^{x/T}]^{-1}$ and $g = 6$. At $T = 0$ the antiparticle
terms vanish, the occupations become step functions at
$k_F = \sqrt{\mu^2 - m^2}$ (with $n_F = 0$ when $\mu \le m$), and
%
\begin{align}
n_F &= \frac{g\, k_F^3}{6\pi^2 \hc} , \\
\varepsilon_F &= \frac{g}{16 \pi^2 \hc}
   \left[ k_F \big( 2 k_F^2 + m^2 \big) E_F
        - m^4 \ln \frac{k_F + E_F}{m} \right] , \\
P_F &= \frac{g}{48 \pi^2 \hc}
   \left[ k_F \big( 2 k_F^2 - 3 m^2 \big) E_F
        + 3 m^4 \ln \frac{k_F + E_F}{m} \right] ,
\qquad s_F = 0 ,
\end{align}
%
with $E_F = \sqrt{k_F^2 + m^2} = \mu$. These integrals are not implemented in
this subpackage: they come from \texttt{eos.general.fermi\_integrals}, which
evaluates them through the Johns--Ellis--Lattimer analytic
approximation~\cite{JohnsEllisLattimer1996}, uniformly valid from the
degenerate to the non-degenerate limit and exact at $T = 0$.

The perturbative correction to a massive flavour is taken to be the same
function of $\mu$ and $T$ as for a massless one:
%
\begin{equation}
X(\mu, T, m, \as) \;=\; X_F(\mu, T, m)
   \;+\; \big[ X_0(\mu, T, \as) - X_0(\mu, T, 0) \big] ,
\qquad X \in \{ n, P, \varepsilon, s \} ,
\label{eq:massive}
\end{equation}
%
where the bracket is $-(2\as/\pi)$ times the $\mu$-dependent part of $X_0$ and
$-(50\as/21\pi)$ times its thermal part. This is a prescription, not an
expansion of the massive result: the true $\mathcal{O}(\as)$ correction to a
gas of mass $m$ differs from Eq.~\eqref{eq:massive} at relative order
$m^2/\mu^2$, which for $m_s = 150$~MeV at $\mu_s \simeq 440$~MeV is a
$\sim 12\%$ correction of a $\sim 19\%$ correction. What the prescription does
guarantee is that (i) it reduces exactly to
Eqs.~\eqref{eq:P0}--\eqref{eq:e0} as $m \to 0$, and (ii) the correction is
itself a consistent set --- $\Delta n = \partial \Delta P/\partial\mu$,
$\Delta s = \partial \Delta P/\partial T$, $\Delta\varepsilon = 3 \Delta P$ ---
so adding it to a free Fermi gas that satisfies Eq.~\eqref{eq:euler} leaves a
sector that still does.

\subsection{Gluons}

Thermal gluons are $g_g = 2 \times 8 = 16$ massless bosons at $\mu = 0$, with
their own correction factor:
%
\begin{equation}
P_g = \frac{8\pi^2}{45}\, \frac{T^4 c_g(\as)}{\hc} , \qquad
\varepsilon_g = 3 P_g , \qquad
s_g = \frac{32\pi^2}{45}\, \frac{T^3 c_g(\as)}{\hc} , \qquad
n_g = 0 .
\label{eq:gluons}
\end{equation}
%
They satisfy $\varepsilon_g + P_g = 4 P_g = T s_g$ and carry no conserved
charge, so switching them on shifts $P$, $\varepsilon$ and $s$ and nothing
else. They vanish identically at $T = 0$. \texttt{gluons} is this model's own
sector flag: no other model in the repository has one.

\section{The colour--flavour locked phase}
\label{sec:cfl}

At high density the favoured ground state of three-flavour quark matter is the
CFL condensate~\cite{AlfordRajagopalWilczek1999}, in which every quark pairs
and the common gap $\Delta$ enters the pressure at order $\Delta^2 \mu^2$. The
package implements this as a second phase, not as a further correction to the
first: it has its own potential, its own solver, and its own closure.

\subsection{The gap}

The gap is not solved for. It is imposed as a BCS-shaped function of
temperature with a single parameter, the zero-temperature gap $\Delta_0$:
%
\begin{equation}
\Delta(T) = \begin{cases}
   \Delta_0 \sqrt{1 - T^2/T_c^2} , & T < T_c , \\[2pt]
   0 , & T \ge T_c ,
\end{cases}
\qquad
T_c = 0.57 \cdot 2^{1/3}\, \Delta_0 = 0.71815\, \Delta_0 ,
\label{eq:gap}
\end{equation}
%
so that $\Delta_0 = 100$~MeV gives $T_c = 71.815$~MeV. The entropy needs the
derivative, which follows by differentiating Eq.~\eqref{eq:gap}:
%
\begin{equation}
\frac{d\Delta}{dT} \;=\; -\,\frac{\Delta_0\, T}{T_c^2 \sqrt{1 - T^2/T_c^2}} ,
\qquad T < T_c ,
\label{eq:dgap}
\end{equation}
%
and zero at $T = 0$ and $T \ge T_c$. Equation~\eqref{eq:dgap} diverges as
$T \to T_c^-$; the code returns zero once the square root falls below
$10^{-10}$, which bounds the entropy correction rather than letting it blow
up at the last grid point before $T_c$.

\subsection{The paired potential}

The pairing term is that of Alford, Braby, Paris and
Reddy~\cite{AlfordBrabyParisReddy2005}, written per flavour so that the three
potentials need not be equal:
%
\begin{equation}
P^{\rm CFL} \;=\; \sum_q P_q(\mu_q, T, m_q, \as)
   \;+\; \frac{\Delta(T)^2}{\pi^2 \hc} \sum_q \mu_q^2
   \;-\; \frac{B}{\hc} .
\label{eq:Pcfl}
\end{equation}
%
When the three potentials coincide at $\bar\mu$ the pairing term is
$3\bar\mu^2\Delta^2/(\pi^2\hc)$, which is the $3\Delta^2\mu^2/\pi^2$ of
Ref.~\cite{AlfordBrabyParisReddy2005} and of \texttt{eos/abpr}. The
$-3 m_s^2 \mu^2/(4\pi^2)$ term those references carry alongside it is
\emph{not} added here: it is the leading expansion of the massive strange
Fermi gas, which Eq.~\eqref{eq:Pcfl} already contains exactly through
$P_s(\mu_s, T, m_s, \as)$, and adding both would count it twice.

Everything else follows from Eq.~\eqref{eq:Pcfl} as derivatives, which is what
makes the paired sector thermodynamically consistent:
%
\begin{align}
n_q &= \frac{\partial P^{\rm CFL}}{\partial \mu_q}
     = n_q(\mu_q, T, m_q, \as)
       + \frac{2 \mu_q \Delta(T)^2}{\pi^2 \hc} ,
\label{eq:ncfl} \\
s &= \frac{\partial P^{\rm CFL}}{\partial T}
   = \sum_q s_q(\mu_q, T, m_q, \as)
     + \frac{2\,\Delta(T)}{\pi^2 \hc} \frac{d\Delta}{dT} \sum_q \mu_q^2 ,
\label{eq:scfl} \\
f &= -P^{\rm CFL} + \sum_q \mu_q n_q , \qquad
\varepsilon = f + T s .
\label{eq:ecfl}
\end{align}
%
The energy density is \emph{defined} by Eq.~\eqref{eq:ecfl} rather than summed
from the flavours, so Eq.~\eqref{eq:euler} holds in the paired phase by
construction; it is checked all the same, and holds to $2\times10^{-16}$
relative. Note that the entropy correction \eqref{eq:scfl} is negative
wherever the gap is falling ($d\Delta/dT < 0$): pairing removes states from
the Fermi surface and the condensate carries less entropy than the gas it
replaces.

The gluon term is not part of Eq.~\eqref{eq:Pcfl}. In the CFL phase the gluons
are all massive through the Meissner effect and their thermal population is
suppressed; the sector remains available as a flag at the solver level and is
added to the totals there, as it is in the unpaired phase, but it is not
inside the phase's own potential.

\section{Conserved charges}
\label{sec:charges}

With $S = +1$ per $s$ quark --- this repository's convention, opposite to the
PDG --- and $C$ the electric charge of the strongly-interacting matter alone,
the quark quantum numbers $(B, C, S)$ are $(\tfrac13, +\tfrac23, 0)$ for $u$,
$(\tfrac13, -\tfrac13, 0)$ for $d$ and $(\tfrac13, -\tfrac13, +1)$ for $s$, so
%
\begin{equation}
\nB = \frac{n_u + n_d + n_s}{3} , \qquad
\nC = \frac{2 n_u - n_d - n_s}{3} , \qquad
\nS = n_s ,
\label{eq:charges}
\end{equation}
%
and the potentials map both ways,
%
\begin{equation}
\muB = \mu_u + 2\mu_d , \qquad
\muC = \mu_u - \mu_d , \qquad
\muS = \mu_s - \mu_d ,
\label{eq:basis}
\end{equation}
%
inverted by $\mu_u = \tfrac13\muB + \tfrac23\muC$,
$\mu_d = \tfrac13\muB - \tfrac13\muC$,
$\mu_s = \tfrac13\muB - \tfrac13\muC + \muS$. These maps are not written in
this subpackage: they are \texttt{eos.general.basis}, shared with every model,
and a test asserts bit-for-bit agreement between the two. The fractions are
$Y_C = \nC/\nB$ and $Y_S = \nS/\nB$, always per baryon and always excluding
the leptons; each flavour also reports $Y_q = n_q/\nB$, so that
$Y_u + Y_d + Y_s = 3$ identically.

The sign of $\muC$ is fixed by $\muC = \mu_u - \mu_d$, which is
$\mu_p - \mu_n$ in the hadronic sector, so beta equilibrium reads
$\muC + \mu_e = 0$ here exactly as it does there.

\section{Leptons, photons and the totals}
\label{sec:totals}

The quantities above are those of the strongly-interacting sector. What a
solved point reports adds, from
\texttt{eos.general.thermodynamics\_leptons}, whichever of these the mode and
the flags call for:
%
\begin{itemize}
\item \textbf{electrons} (with positrons) at $\mu_e$, a Fermi gas of mass
      $0.511$~MeV and degeneracy $2$, present wherever the mode has a lepton
      condition;
\item \textbf{electron neutrinos} (with antineutrinos) at $\munue$, massless
      and of degeneracy $1$, present only in the trapped mode;
\item \textbf{photons}, $P_\gamma = \tfrac13\varepsilon_\gamma =
      \pi^2 T^4/(45\hc)$, $s_\gamma = 4\varepsilon_\gamma/(3T)$;
\item \textbf{thermal neutrinos}: the flavours \emph{not} tracked in the
      composition, carried as $\mu = 0$ gases. Three flavours where the
      electron neutrino is free-streaming ($\munue = 0$), two where it is
      trapped and therefore already counted at its own potential.
\end{itemize}
%
None of these carries $B$, $C$ or $S$ in the sense of Eq.~\eqref{eq:charges} ---
$\nC$ is non-leptonic by definition --- so they enter $P$, $\varepsilon$ and
$s$ and the neutrality condition, and nothing else. The totals are
%
\begin{equation}
P = P^{\rm matter} + P_e + P_{\nu_e} + P_\gamma + P_{\nu}^{\rm th} , \qquad
\varepsilon = \varepsilon^{\rm matter} + \varepsilon_e + \varepsilon_{\nu_e}
              + \varepsilon_\gamma + \varepsilon_{\nu}^{\rm th} , \qquad
s = s^{\rm matter} + s_e + s_{\nu_e} + s_\gamma + s_{\nu}^{\rm th} ,
\end{equation}
%
with $P^{\rm matter}$ from Eq.~\eqref{eq:P} or Eq.~\eqref{eq:Pcfl} and
$f = \varepsilon - Ts$.

\section{Equilibrium modes and their closures}
\label{sec:modes}

A mode fixes which quantities are imposed and which are unknown. In the
unpaired phase the unknowns are chemical potentials only --- the potential of
Sec.~\ref{sec:potential} is explicit in $\mu$, so unlike a mean-field model
there is no field equation to carry along and no density in the unknown
vector.

\paragraph{\texttt{beta\_eq\_neutrinoless}} $(\nB, T)$; unknowns
$x = [\mu_u, \mu_d, \mu_s, \mu_e]$, four rows:
%
\begin{equation}
\begin{aligned}
r_1 &= \frac{n_u + n_d + n_s}{3} - \nB ,
&\qquad r_2 &= \nC - n_e(\mu_e, T) , \\
r_3 &= \mu_d - \mu_u - \mu_e \;\; (= -\muC - \mu_e) ,
&\qquad r_4 &= \mu_s - \mu_d \;\; (= \muS) .
\end{aligned}
\label{eq:res_beta}
\end{equation}
%
Row $r_1$ is baryon number, $r_2$ total electric neutrality, $r_3$ beta
equilibrium $d \leftrightarrow u + e^- + \bar\nu_e$ with free-streaming
neutrinos, and $r_4$ strangeness equilibrium $s \leftrightarrow d$, which is
the statement $\muS = 0$.

\paragraph{\texttt{beta\_eq\_neutrino\_trapped}} $(\nB, Y_{L_e}, T)$; unknowns
$x = [\mu_u, \mu_d, \mu_s, \mu_e, \munue]$, five rows: $r_1$, $r_2$ and $r_4$
unchanged, with
%
\begin{equation}
r_3 = \mu_d - \mu_u - \mu_e + \munue \;\; (= -\muC - \mu_e + \munue) ,
\qquad
r_5 = \frac{n_e + n_{\nu_e}}{\nB} - Y_{L_e} .
\end{equation}
%
The neutrino gas carries its antiparticles, so $n_{\nu_e}$ is the net density
and $Y_{L_e}$ is the conserved electron-family number per baryon. The muon
family is not tracked, so $Y_{L_\mu}$ raises.

\paragraph{\texttt{fixed\_YC}} $(\nB, Y_C, T)$; unknowns
$x = [\mu_u, \mu_d, \mu_s]$, three rows: $r_1$ and $r_4$ unchanged, with the
charge fraction imposed instead of neutrality,
%
\begin{equation}
r_2' = \frac{\nC}{\nB^{\rm calc}} - Y_C , \qquad
\nB^{\rm calc} = \frac{n_u + n_d + n_s}{3} .
\end{equation}
%
With \texttt{leptons=False} the result is electrically charged quark matter,
which is what a mixed-phase construction needs per pure phase before global
neutrality is imposed. With \texttt{leptons=True} a neutralizing electron gas
is added \emph{after} the solve, by inverting $n_e(\mu_e, T) = \nC$ for
$\mu_e$ --- a one-dimensional inversion, not a further row, because the
strongly-interacting sector of this model does not respond to $\mu_e$ at fixed
$Y_C$. The two therefore share the same quark state and differ only in the
lepton contribution to $P$, $\varepsilon$ and $s$.

\paragraph{\texttt{fixed\_YC\_YS}} $(\nB, Y_C, Y_S, T)$; unknowns
$x = [\mu_u, \mu_d, \mu_s]$, three rows: $r_1$ and $r_2'$ unchanged, with
strangeness equilibrium replaced by the imposed strangeness fraction,
%
\begin{equation}
r_3' = \frac{n_s}{\nB^{\rm calc}} - Y_S .
\end{equation}
%
This is the mode that separates a bag model from a nucleonic one: it is
meaningful here, where the $s$ quark is a degree of freedom, and it raises in
\texttt{eos/zl}, which has none. Leptons are handled as in
\texttt{fixed\_YC}.

\paragraph{The CFL phase} is a phase selector, not one of the four modes. It
is closed by flavour locking rather than by an equilibrium condition: the
condensate pairs the three flavours in equal numbers, so
%
\begin{equation}
r_q = n_q^{\rm CFL}(\mu_q, T, m_q, \as, \Delta_0) - \nB = 0 , \qquad q = u, d, s ,
\label{eq:res_cfl}
\end{equation}
%
with $n_q^{\rm CFL}$ from Eq.~\eqref{eq:ncfl}, is three rows for the three
unknowns $x = [\mu_u, \mu_d, \mu_s]$ at given $(\nB, T, \Delta_0)$. Two
consequences are worth stating because they are easy to expect wrongly:

\begin{itemize}
\item $\nC = (2 n_u - n_d - n_s)/3 = 0$ identically, so the phase is
      \emph{electrically neutral by construction} with no electrons at all,
      and $Y_C$ is returned as zero to round-off ($-1.7\times10^{-14}$ at
      $\nB = 0.8\ \mathrm{fm}^{-3}$) rather than solved for. This is the
      physical content of CFL neutrality~\cite{AlfordRajagopalWilczek1999}.
\item The phase is \emph{not} in strangeness equilibrium. Equal densities at
      unequal masses require unequal potentials --- at $\nB = 0.8\
      \mathrm{fm}^{-3}$, $T = 0$, $\Delta_0 = 100$~MeV the solve gives
      $\mu_u = \mu_d = 402.168$~MeV and $\mu_s = 434.121$~MeV --- so
      $\muS = \mu_s - \mu_d \neq 0$, and $\varepsilon/\nB = \muB + \muS$ at
      $P = 0$ rather than $\muB$. Reading $\muB$ as the energy per baryon of
      the paired phase is the standard way to misread it.
\end{itemize}

\medskip
Wherever a temperature is accepted, entropy per baryon may be given in its
place, through an outer one-dimensional solve for $T$ at fixed $\nB$ and
fractions.

\section{What a solved point returns}
\label{sec:returns}

Every mode returns the same record, so a caller reads one shape whichever mode
it asked for:

\begin{center}
\begin{tabular}{lll}
\toprule
field & symbol & note \\
\midrule
\texttt{converged}, \texttt{error} & --- & the status and the largest scaled
                                            residual (Sec.~\ref{sec:numerics}) \\
\texttt{n\_B}, \texttt{T} & $\nB$, $T$ & the conditions the point was solved at \\
\texttt{Y\_C}, \texttt{Y\_S}, \texttt{Y\_L} & $Y_C, Y_S, Y_{L_e}$ & as
                                            realised, Eq.~\eqref{eq:charges} \\
\texttt{mu\_u}, \texttt{mu\_d}, \texttt{mu\_s} & $\mu_u, \mu_d, \mu_s$ & the unknowns \\
\texttt{mu\_e}, \texttt{mu\_nu} & $\mu_e, \munue$ & zero where the mode has no
                                            lepton condition \\
\texttt{mu\_B}, \texttt{mu\_C}, \texttt{mu\_S} & $\muB, \muC, \muS$ & derived,
                                            Eq.~\eqref{eq:basis} \\
\texttt{n\_u}, \texttt{n\_d}, \texttt{n\_s} & $n_u, n_d, n_s$ & net densities,
                                            antiquarks subtracted \\
\texttt{n\_e}, \texttt{n\_nu} & $n_e, n_{\nu_e}$ & net, antiparticles subtracted \\
\texttt{P\_total}, \texttt{e\_total} & $P$, $\varepsilon$ & MeV/fm$^3$,
                                            Sec.~\ref{sec:totals} \\
\texttt{s\_total}, \texttt{f\_total} & $s$, $f = \varepsilon - Ts$ & fm$^{-3}$,
                                            MeV/fm$^3$ \\
\texttt{Y\_u}, \texttt{Y\_d}, \texttt{Y\_s}, \texttt{Y\_e}, \texttt{Y\_nu}
      & $n_i/\nB$ & per baryon \\
\bottomrule
\end{tabular}
\end{center}
%
The CFL record carries the same fields and adds \texttt{Delta0} and
\texttt{Delta} --- the parameter and its value at $T$, Eq.~\eqref{eq:gap} ---
and \emph{drops} \texttt{n\_e} and \texttt{n\_nu}, since the phase has no
leptons; \texttt{mu\_e}, \texttt{mu\_nu}, \texttt{Y\_e} and \texttt{Y\_nu}
are present and stay at zero. (\texttt{CFLPoint}'s own docstring says
\texttt{Y\_e} is absent as well; it is not.)

There is no scalar density in this model and therefore no identity
$n_s = (\varepsilon - 3P)/m^*$ to state: the quark masses are parameters and
there is no gap equation, so nothing here plays the part $m^*$ plays in a
mean-field model. The \texttt{n\_s} of the returns table is the strange-quark
number density.

The entropy density is worth a word because nothing in the residual derives
from it and it is therefore easy to leave unstated. It is not integrated
separately: for a massless flavour it is the closed form
Eq.~\eqref{eq:s0}, for a massive one it is Eq.~\eqref{eq:sF} --- the Euler
relation of the free gas, which is exact --- plus the correction of
Eq.~\eqref{eq:massive}, and in the paired phase it carries the further term of
Eq.~\eqref{eq:scfl}. The gluon, photon and lepton entropies are added on top.

\section{The API surface}
\label{sec:api}

Three entry points, with the signatures every model in this repository carries:
%
\begin{verbatim}
eos_point(par, mode, species, n_B=, T=|SnB=, leptons=, x0=, **conditions)
eos_table(par, mode, species, axes=, fixed=, leptons=, skip_errors=,
          rows=, progress=, verbose=)
eos_response(par, mode, species, frozen="equilibrium", n_B=, T=,
             leptons=, rel_step=1e-3, **conditions)
\end{verbatim}
%
\texttt{par} comes first and is never optional. \texttt{eos\_point} returns a
\texttt{PointResult} whose \texttt{.ok} must be tested before \texttt{.point}
is read; non-convergence is a return value, not an exception.
\texttt{eos\_table} takes \texttt{axes=\{'nB': grid, 'T': grid\}} plus the
mode's own fraction axis (\texttt{Y\_C}, \texttt{Y\_S}, \texttt{Y\_Le}, or
\texttt{Delta0} for the paired phase) and returns arrays consumable by
\texttt{eos.astro.tov} with no adapter; its \texttt{progress} callback is
invoked once per completed line with the repository's standard dictionary
\texttt{\{mode, line, n\_lines, temp\_key, temp, fracs, n\_solved,
n\_requested, elapsed\_s\}}, and \texttt{verbose=True} installs the built-in
printer. Deep solver code never prints.

\texttt{eos\_response} implements one freeze,
\texttt{frozen="equilibrium"}: everything re-equilibrates under the
perturbation, so the derivatives are taken along the mode's own sequence,
%
\begin{equation}
c_{s,\rm eq}^2 = \left. \frac{\partial P}{\partial \varepsilon} \right|_{T} ,
\qquad
C_V = \frac{T}{\nB} \left. \frac{\partial s}{\partial T} \right|_{\nB} ,
\label{eq:response}
\end{equation}
%
both by central differences over a relative step \texttt{rel\_step} in the
variable differentiated, since the residual of this model has no analytic
Jacobian in this repository. $C_V$ is returned only at $T > 0$, where it is
defined. The remaining freezes of the specification --- frozen per-species
composition, frozen conserved fractions, the leptonic re-neutralization
variants --- raise \texttt{NotImplementedError} naming the gap, and are
recorded in \texttt{docs/DEFERRED.md}. With
\texttt{frozen="equilibrium"} there is no isothermal/adiabatic ambiguity left
to name: Eq.~\eqref{eq:response} is taken at fixed $T$ and its name says so.

\section{Numerics}
\label{sec:numerics}

Each mode is a system of three to five equations, solved with Powell's hybrid
method and, if that does not reach the gate, with Levenberg--Marquardt; when a
warm start was supplied and both fail, one further hybrid attempt is made from
the mode's own cold guess, since a warm start carried across the strange-quark
onset can land outside the basin. Every attempt is bounded internally and at
most three are made: a parameter scan must always get an answer back, and
non-convergence is a value carried on the record, never an exception and never
an unbounded loop.

Convergence is judged on a dimensionless residual. The rows carry mixed units
--- densities of order $10^{-1}\,\mathrm{fm}^{-3}$, fractions of order unity,
equalities between chemical potentials of order $10^{3}$~MeV --- so each row
is divided by the scale of the quantity it balances: $\nB$ for a density row
($r_1$, $r_2$, and the three rows of the CFL system), $|\muB|$ for a potential
equality ($r_3$, $r_4$), and unity for a row that is already dimensionless
($r_2'$, $r_3'$, $r_5$). The state is accepted when the largest scaled
component is below $10^{-10}$. A gate on the raw residual vector is dominated
by whichever row happens to be largest and accepts states satisfying the
others only loosely; a gate on the solver's own success flag reports whether
the iteration terminated, which is a different question again.

The cold guess estimates $\mu \simeq (\pi^2 \nB)^{1/3} \hbar c$ from the
massless relation $n \simeq \mu^3/(\pi^2\hc)$ at one flavour per baryon, and
adjusts it per mode. Tables are built line by line --- one line per
temperature and per combination of the fractions the mode fixes --- and swept
along the baryon density with a warm start, each solved point seeding the
next, with the step bisected where a solve misses so that the sweep can walk
through the strange-quark onset rather than stopping at it. That loop is
\texttt{eos.general.tabulate}, shared with the other models; what this
subpackage supplies is which solver a mode name means and what part of a
solved point becomes the next guess.

\begin{acknowledgments}
Written as the specification of \texttt{eos/alphabag}.
\end{acknowledgments}

\bibliography{../../docs/eos}

\end{document}
